\section{Carrier splitting under a new Boolean observation}\label{sec:carrier-splitting}

We now treat the refinement regime left open in the previous draft: a new
observation is not constant on every current observational-equivalence class, so
one old $T_0$ state can split into two refined states.  This operation is not a
cover refinement and is not, on its face, a same-carrier order thinning.  It does,
however, admit a canonical factorization through one.

Let $W$ be the finite world set, let $\Phi$ be the current observation family,
and write
\[
 Q=o_\Phi(W)\subseteq\{0,1\}^{I}
\]
for the finite $T_0$ quotient ordered coordinatewise.  Add one Boolean
observation $\psi$.  The refined quotient is
\[
 Q^\psi=o_{\Phi\cup\{\psi\}}(W)
 \subseteq Q\times\{0,1\},
\]
with the product order.  Let
\[
 \pi:Q^\psi\longrightarrow Q,
 \qquad \pi(s,a)=s,
\]
be the coordinate-forgetting map.  For $a\in\{0,1\}$ put
\[
 A_a:=\{s\in Q:(s,a)\in Q^\psi\},
 \qquad C:=A_0\cap A_1.
\]
Thus $A_0\cup A_1=Q$, and $C$ is precisely the set of old observational
states that split.

\begin{definition}[Split completion]
On the same carrier as $Q^\psi$, define the split-completion order
$\widehat Q^\psi$ by
\[
 (s,a)\leq_{\widehat Q}(t,b)
 \quad\Longleftrightarrow\quad
 \bigl(s<t\text{ in }Q\bigr)
 \ \text{or}\ 
 \bigl(s=t\text{ and }a\leq b\bigr).
\]
Equivalently, $\widehat Q^\psi$ is the lexicographic sum over $Q$ of the
nonempty fiber chains
\[
 F_s:=\{a\in\{0,1\}:(s,a)\in Q^\psi\}.
\]
Let $b(s,a)=a$ be the new-bit label on this expanded carrier.
\end{definition}

The distinction between $Q^\psi$ and $\widehat Q^\psi$ is exact.  The split
completion remembers every old comparison between coarse states, independently
of the new bit.  The actual refined order keeps only those comparisons along
which the new bit is nondecreasing.

\begin{theorem}[Carrier-splitting factorization]\label{thm:split-factorization}
With the notation above:
\begin{enumerate}[label=(\roman*)]
\item the projection $p:\widehat Q^\psi\to Q$, $p(s,a)=s$, is a homotopy
  equivalence of finite posets; in fact $\widehat Q^\psi$ beat-reduces to a
  subposet isomorphic to $Q$;
\item the actual refined quotient is exactly the one-bit thinning of the split
  completion,
  \[
    Q^\psi=(\widehat Q^\psi)_b,
  \]
  meaning that a comparison $x<_{\widehat Q}y$ is retained precisely when
  $b(x)\leq b(y)$;
\item if $i:\Delta(Q^\psi)\hookrightarrow\Delta(\widehat Q^\psi)$ is the
  inclusion, then $\Delta(\pi)=\Delta(p)\circ i$, and there is a homotopy
  equivalence of homotopy cofibers
  \[
  \operatorname{Cof}(\Delta\pi)\simeq
  \operatorname{Cof}(i).
  \]
  Consequently, for every coefficient ring $R$,
  \[
  \widetilde H_n\!\left(\Cone(\Delta\pi);R\right)
  \cong
  H_n\!\left(\Delta(\widehat Q^\psi),\Delta(Q^\psi);R\right).
  \]
\end{enumerate}
\end{theorem}

\begin{proof}
For every split state $s\in C$, the lower copy $(s,0)$ is an up-beat point of
$\widehat Q^\psi$: its strict upper set has the minimum $(s,1)$, because
$(s,1)$ lies below every copy over every $t>s$.  Remove these lower split
copies one at a time.  What remains has one representative over each $s\in Q$
and is order-isomorphic to $Q$.  Beat-point removal is a strong deformation
retraction for finite $T_0$ spaces, giving (i).

For distinct coarse states $s<t$, the split completion always contains
$(s,a)<(t,b)$ whenever those copies exist.  The product order on $Q^\psi$
retains exactly the cases $a\leq b$.  The same statement holds within a split
fiber, where $(s,0)<(s,1)$ is retained.  This proves (ii).  Part (iii) follows
from $\pi=p\circ i$ and homotopy invariance of the homotopy cofiber under
postcomposition by the homotopy equivalence $\Delta(p)$.  Since $i$ is a
simplicial inclusion, its homotopy cofiber computes the relative homology of
its simplicial pair.
\end{proof}

\begin{remark}[What is standard and what is specific]
Beat-point reduction is classical finite-space topology, and homotopy cofiber
invariance is standard.  The project-specific content is the observation-theoretic
factorization: an actual split of $T_0$ observational states is converted canonically
into a homotopy-neutral lexicographic expansion followed by the precise one-bit
monotonicity thinning already studied on a fixed carrier.  Thus carrier splitting
is not a fourth unrelated operation; it is a neutral blow-up plus an order defect.
\end{remark}

\subsection{The exact inversion-chain complex}

Theorem~\ref{thm:split-factorization} gives a chain-level classification without
any nerve hypothesis.

\begin{theorem}[Inversion-chain model]\label{thm:inversion-chain}
For any coefficient ring $R$, the relative group
\[
 C_k\!\left(\Delta(\widehat Q^\psi),\Delta(Q^\psi);R\right)
\]
is free on the strict chains
\[
 x_0<_{\widehat Q}x_1<_{\widehat Q}\cdots<_{\widehat Q}x_k
\]
whose bit word
\[
 b(x_0)b(x_1)\cdots b(x_k)
\]
is not weakly increasing.  Equivalently, a relative generator is exactly a
split-completion chain containing an inversion $1\cdots0$.  Its relative boundary
is the ordinary alternating simplicial boundary with every monotone-bit face set
to zero.
\end{theorem}

\begin{proof}
The two order complexes have the same vertex set.  A split-completion chain is a
simplex of $\Delta(Q^\psi)$ exactly when every comparison in the chain survives
the one-bit thinning, which is equivalent to the bit word being weakly increasing.
Quotienting the simplicial chain groups therefore leaves exactly the nonmonotone
chains, and the relative differential is the induced simplicial differential.
\end{proof}

This formulation is useful computationally: the defect depends only on the old
signature poset and the three-valued fiber profile
\[
 F_s\in\bigl\{\{0\},\{1\},\{0,1\}\bigr\},
\]
not on the number of worlds inside each old equivalence class.

\subsection{Immediate homotopy-neutrality criteria}

The factorization also reveals when splitting is automatically conservative.

\begin{proposition}[Monotone-section criterion]\label{prop:monotone-section}
If $A_1$ is an upward-closed subset of $Q$, then $\pi:Q^\psi\to Q$ is a
homotopy equivalence.  Dually, if $A_0$ is downward closed, then $\pi$ is a
homotopy equivalence.
\end{proposition}

\begin{proof}
Assume $A_1$ is upward closed.  Define
\[
 r_+(s)=
 \begin{cases}
 (s,1),&s\in A_1,\\
 (s,0),&s\notin A_1.
 \end{cases}
\]
The upset condition makes $r_+$ order preserving, and $\pi r_+=\mathrm{id}_Q$.
Moreover every $(s,a)\in Q^\psi$ satisfies $(s,a)\leq r_+(s)$, so
$\mathrm{id}_{Q^\psi}\leq r_+\pi$.  The order-homotopy lemma gives
$r_+\pi\simeq\mathrm{id}$ and hence $\pi$ is a homotopy equivalence.  The
$A_0$ statement is dual, using the minimum available bit in each fiber.
\end{proof}

\begin{corollary}[Common-outcome anchor]
If every old observational state has at least one refined world with $\psi=1$
($A_1=Q$), or every old state has at least one with $\psi=0$ ($A_0=Q$), then
carrier splitting is homotopy-neutral.  In particular, if every old state genuinely
splits, so $A_0=A_1=Q$, then $Q^\psi=Q\times\{0<1\}$ and the projection is a
homotopy equivalence.
\end{corollary}

Thus the act of duplicating a state is not itself the source of topological defect.
The defect is caused by incompatible outcome availability across comparable coarse
states: a lower state that can realize $1$ together with an upper state that can
realize $0$ creates a potentially deleted comparison.

\subsection{Localization by lower fibers}

For $s\in Q$ define the lower fiber
\[
 L_s:=\pi^{-1}(Q_{\leq s}).
\]
If $s\in A_1$, then $(s,1)$ is a maximum of $L_s$, hence $L_s$ is
contractible.  Therefore all lower-fiber obstructions to Quillen-type equivalence
localize to the $0$-only states $A_0\setminus A_1$.  Dually, all upper-fiber
obstructions localize to the $1$-only states $A_1\setminus A_0$.

This observation does not replace the exact relative complex above.  It does,
however, connect the carrier-splitting problem to the classical poset-fiber
framework.  Under the connectivity hypotheses of the Bj\"orner--Wachs--Welker
fiber theorem, the source $\Delta(Q^\psi)$ admits a wedge decomposition in terms
of the noncontractible lower fibers and upper links.  The important point here is
that the split fibers themselves disappear from that list whenever their top copy
is present.

\subsection{Exact low-height classification}

The general inversion complex can have arbitrarily high dimension.  In the first
nontrivial dimensional regime it collapses to graph incidence and admits a closed
integral classification.

\begin{definition}[Weighted split height]
Let $C=A_0\cap A_1$ be the split set.  Define
\[
 h_\psi(Q)
 :=\max_{s_0<\cdots<s_r\text{ in }Q}
 \left(r+\#\{j:s_j\in C\}\right).
\]
\end{definition}

\begin{lemma}\label{lem:split-height}
$h_\psi(Q)=\dim\Delta(\widehat Q^\psi)$.
\end{lemma}
\begin{proof}
A strict coarse chain contributes one vertex from every fiber, and a split fiber can
contribute both copies consecutively.  Taking both copies at every split state on the
chain attains the displayed number, and no split-completion chain can contain more.
\end{proof}

Suppose $h_\psi(Q)\leq2$.  An \emph{inversion edge} is a deleted order-complex edge
of $\Delta(\widehat Q^\psi)$, equivalently a pair
\[
 (s,1)<_{\widehat Q}(t,0),\qquad s<t,
\]
with $s\in A_1$ and $t\in A_0$.  A deleted triangle has one of the four nonmonotone
bit patterns
\[
 010,\qquad100,\qquad101,\qquad110.
\]
It therefore contains either one or two inversion edges.  Introduce one formal root
for each connected component of $Q$ (equivalently of $\widehat Q^\psi$).  The
\emph{carrier-split defect graph} $G_\psi(Q)$ has one vertex for every inversion edge;
a deleted triangle with two inversion edges joins their vertices, while a deleted
triangle with one inversion edge joins that vertex to the root of its component.
Parallel graph edges are allowed.

\begin{theorem}[Carrier-splitting cone classification in weighted height two]
\label{thm:carrier-graph}
Assume $h_\psi(Q)\leq2$.  Let $c_0(G_\psi)$ be the number of connected components
of $G_\psi(Q)$ containing no formal root.  Then
\[
 \widetilde H_n\!\left(\Cone(\Delta\pi);\mathbb Z\right)
 \cong
 \begin{cases}
   \mathbb Z^{\beta_1(G_\psi)},&n=2,\\
   \mathbb Z^{c_0(G_\psi)},&n=1,\\
   0,&\text{otherwise}.
 \end{cases}
\]
The relative boundary matrix is totally unimodular.  Hence the same Betti ranks hold
over every field.  Moreover the following are equivalent:
\begin{enumerate}[label=(\roman*)]
\item the carrier-splitting cone is acyclic over $\mathbb Z$;
\item it is acyclic over every field;
\item $G_\psi(Q)$ is a rooted forest;
\item $\Delta(\widehat Q^\psi)$ collapses onto $\Delta(Q^\psi)$ through deleted
edge--triangle pairs;
\item $\Delta\pi:\Delta(Q^\psi)\to\Delta(Q)$ is a homotopy equivalence.
\end{enumerate}
\end{theorem}

\begin{proof}
By Theorem~\ref{thm:split-factorization}, the mapping-cone homology is the relative
homology of the one-bit thinning
$\Delta(Q^\psi)\subseteq\Delta(\widehat Q^\psi)$.  Lemma~\ref{lem:split-height}
places the split completion in dimension at most two.  The one-bit height-two
incidence theorem from the companion order-thinning analysis applies verbatim: the
relative $C_2\to C_1$ boundary is the oriented incidence matrix of the rooted defect
graph with root rows deleted.  Ordinary graph incidence gives the two displayed free
homology groups and total unimodularity; the rooted-forest condition is equivalent to
perfect relative cancellation and to an elementary collapse.  Finally
$\Delta(p)$ is already a homotopy equivalence, so two-out-of-three identifies
homotopy neutrality of the inclusion with homotopy neutrality of $\Delta\pi$.
\end{proof}

\begin{remark}[Status of the low-height theorem]
The graph-incidence algebra, total unimodularity, rooted forests, and beat-point
technology are standard.  The exact height-two theorem for one-bit thinning is proved
in the companion structural-computation manuscript.  The new point here is the
carrier-splitting reduction that makes that theorem applicable to a process which
changes the $T_0$ carrier.  Accordingly, this theorem should be positioned as an
exact transfer theorem, not as a new general theory of rooted forests.
\end{remark}

\subsection{Sharpness beyond weighted height two}

No coefficient-independent forest criterion can hold in unrestricted dimension.
The carrier-splitting family contains the no-split case $C=\varnothing$, which is
exactly arbitrary same-carrier one-bit thinning.  Companion results show that one-bit
thinning in height three can already have relative integral torsion, and in unrestricted
height can encode arbitrary finite-poset homotopy types inside a contractible coarse
cone.  If one insists that at least one old state genuinely split, a disjoint neutral
split component may be adjoined without changing positive-degree cone homology.
Thus the weighted-height-two hypothesis in Theorem~\ref{thm:carrier-graph} marks a
real boundary rather than a presentational convenience.

The appropriate all-height object is therefore the inversion-chain complex of
Theorem~\ref{thm:inversion-chain}, supplemented where useful by poset-fiber or
homotopy-colimit methods.  A realistic next target is not another universal forest
theorem, but a classification of higher-dimensional inversion complexes under
additional structural hypotheses on $Q$ and on the split profile $(A_0,A_1)$.
